\documentclass[11pt,a4paper]{article}
\usepackage[T1]{fontenc}
\usepackage[utf8]{inputenc}
\usepackage{lmodern}
\usepackage[margin=26mm,headheight=15pt]{geometry}
\usepackage{amsmath,amssymb,amsthm,mathtools}
\usepackage{microtype,booktabs,array,longtable}
\usepackage{enumitem,fancyhdr,xcolor,xurl}
\usepackage[colorlinks=true,linkcolor=blue!40!black,citecolor=blue!40!black,urlcolor=blue!40!black]{hyperref}
\hypersetup{pdftitle={At the Critical Line: Continued Fractions and Sharp Riesz Summation of Resonant Transseries},pdfauthor={Research report prepared for Vladimir Reshetnikov}}
\setlength{\emergencystretch}{2em}
\setlist{itemsep=3pt,topsep=5pt}
\numberwithin{equation}{section}
\newtheorem{theorem}{Theorem}[section]
\newtheorem{lemma}[theorem]{Lemma}
\newtheorem{proposition}[theorem]{Proposition}
\newtheorem{corollary}[theorem]{Corollary}
\theoremstyle{definition}
\newtheorem{definition}[theorem]{Definition}
\newtheorem{example}[theorem]{Example}
\theoremstyle{remark}
\newtheorem{remark}[theorem]{Remark}
\newcommand{\N}{\mathbb N}
\newcommand{\Z}{\mathbb Z}
\newcommand{\R}{\mathbb R}
\newcommand{\C}{\mathbb C}
\newcommand{\e}{\mathrm e}
\newcommand{\ii}{\mathrm i}
\newcommand{\dd}{\,\mathrm d}
\newcommand{\Res}{\operatorname{Res}}
\newcommand{\re}{\operatorname{Re}}
\newcommand{\im}{\operatorname{Im}}
\newcommand{\dist}{\operatorname{dist}}
\newcommand{\Hh}{\mathbb H}
\newcommand{\Ssum}{\mathcal S}
\newcommand{\Riesz}{\mathcal R}
\newcommand{\Accel}{\mathcal A}
\newcommand{\norm}[1]{\left\lVert #1\right\rVert}
\newcommand{\repo}{\textit{ProveIt}}
\pagestyle{fancy}
\fancyhf{}
\fancyhead[L]{\small At the Critical Line}
\fancyhead[R]{\small Resonant transseries and summation}
\fancyfoot[C]{\thepage}

\begin{document}
\hypersetup{pageanchor=false}
\begin{titlepage}
\vspace*{8mm}
\noindent{\large\scshape Transseries / Small Divisors / Summability}\par
\vspace{12mm}
\noindent{\Huge\bfseries At the Critical Line}\par
\vspace{6mm}
\noindent{\LARGE Continued fractions and sharp Riesz\par\noindent summation of resonant transseries}\par
\vspace{11mm}
\noindent Research article prepared for Vladimir Reshetnikov\par
\noindent 29 September 2026\par
\vspace{10mm}
\begin{abstract}
The residue expansion of a meromorphic Borel transform can diverge even
when its Stokes jump is holomorphic and has a normally convergent
finite-block representation. We resolve the critical-line question
explicitly left open in the \repo\ resonance-block report for two sine
lattices with numerator $t^D$, $D\ge2$. If $p_k/q_k$ are the convergents of an irrational period ratio,
$b=\limsup\log q_{k+1}/q_k\in(0,\infty)$, and
$A_k=q_k^Dq_{k+1}\e^{-bq_k}$, then the increasing-action residue series
converges on the entire line $\re z=b$ exactly when $A_k\to0$, and
converges absolutely exactly when $\sum A_k<\infty$. The two separately
indexed lattice series have a different exact criterion: convergence of
$\sum(-1)^{p_k+q_k+k}A_k\e^{-\ii q_k\im z}$. The proof removes an
absolutely summable remainder and isolates the primitive convergent
pairs. Continued-fraction constructions realize absolute convergence,
conditional convergence, bounded divergent spikes, and unbounded spikes
at every prescribed positive critical abscissa. A derivative hierarchy
shows precisely when an analytic sum permits termwise differentiation of
its raw representation.

For $d$ sine lattices we also prove that Riesz weighting of integer order
$m\ge d-1$ recovers the block sum throughout $\re z>0$, without any
Diophantine hypothesis. It has an exact finite derivative bias and an
exponentially small remainder, uniform across arbitrarily close poles.
The threshold $d-1$ is sharp for gap-independent uniform control, while
continuity at a collision exactly on the cutoff requires one additional
order. Richardson extrapolation removes the bias, and a local inverse
stability theorem transports the exponential accuracy through nonlinear
inversion. All new assertions are proved conventionally; 192 exact
finite checks and 49 high-precision diagnostics accompany the article.
No Lean verification or independently established publication priority
is claimed.
\end{abstract}
\vfill
\noindent\small Inspected repository snapshot:\par
\noindent\small\texttt{3d5973524506411392a911470b5ddc35521568ea}.\par
\noindent\small The source, PDF, verification code, recorded data, and build instructions
are distributed together. No repository files were modified.
\end{titlepage}
\hypersetup{pageanchor=true}
\tableofcontents
\newpage

\section{The open question and the results}
\label{sec:introduction}

\subsection{The precise repository frontier}
The \repo\ report \emph{Resonance-Block Summation of Transseries}
\cite{repo-block} studies
\begin{equation}
 \widehat F_{\boldsymbol\omega}(t)
 =\frac{t^d}{\prod_{j=1}^{d}\sin(\pi\omega_jt)},\qquad \omega_j>0.
 \label{eq:original-kernel}
\end{equation}
It proves a separation-independent block representation of the Stokes
jump, identifies the open half-plane of absolute convergence for two
irrationally related lattices, and transports the block sum through
analytic inversion. These are antecedents, not results claimed here for
the first time.

Its further-research Question~1 asks for a continued-fraction criterion
on the boundary of that half-plane, distinguishing separate lattice
summation from increasing-action summation and constructing the possible
convergence types. The main result below answers that question. The
comparison is to the pinned source, especially its labels
\texttt{thm:block}, \texttt{thm:arithmetic}, and Section~\texttt{sec:questions},
Question~1. The repository inventory \cite{repo-index} was also checked;
this is not an exhaustive priority audit of every document in the
repository or in the literature.

There is a useful distinction behind the problem. A complete close pair
has a bounded sum, but either constituent can be enormous. A cutoff
between the two poles reveals one constituent. Whether those unfinished
pairs vanish is a different question from whether all their absolute
values are summable. Continued fractions identify exactly which pairs
matter at the critical line.

\subsection{An overview in the original degree-two model}
For $(\omega_1,\omega_2)=(1,\alpha)$ with $\alpha>0$ irrational, put
\begin{equation}
 b=\beta(\alpha)=\limsup_{n\to\infty}\frac{1}{n}
                  \log\frac{1}{\dist(n\alpha,\Z)}.
 \label{eq:intro-beta}
\end{equation}
The predecessor proves absolute convergence for $\re z>b$ and failure of
the term test for $\re z<b$. For $0<b<\infty$, let $p_k/q_k$ be the
continued-fraction convergents and set
\begin{equation}
 A_k=q_k^2q_{k+1}\e^{-bq_k},\qquad
 \eta_k=(-1)^{p_k+q_k+k}.
 \label{eq:intro-amplitudes}
\end{equation}
Our increasing-action classification on $\re z=b$ is
\begin{center}
\begin{tabular}{p{.42\textwidth}p{.46\textwidth}}
\toprule
Condition on $A_k$ & Raw residue series, ordered by action\\
\midrule
$\sum_k A_k<\infty$ & Absolutely convergent.\\
$A_k\to0$ and $\sum_k A_k=\infty$ & Conditionally convergent.\\
$A_k\not\to0$ & Divergent; its terms do not tend to zero.\\
\bottomrule
\end{tabular}
\end{center}
The answer is the same at every imaginary height. By contrast, the two
separate lattice series converge at $b+\ii y$ precisely when
$\sum_k\eta_k A_k\e^{-\ii yq_k}$ converges. In particular, an
increasing-action sum can converge although its separate lattice sums
are respectively $+\infty$ and $-\infty$ at the real boundary point.

The second part of the article supplies a different repair when ordinary
summation fails. Write $\lambda$ for an action and $c_\lambda$ for its
Borel residue. For two lattices,
\begin{equation}
 \Riesz_T^{(1)}(z)
 =\sum_{\lambda<T}\left(1-\frac{\lambda}{T}\right)c_\lambda\e^{-\lambda z}
 =\Ssum(z)+\frac{\Ssum'(z)}{T}+O_K(T\e^{-\sigma T}),
 \label{eq:intro-riesz}
\end{equation}
where $K\Subset\{\re z>0\}$ and $\sigma=\min_K\re z$. Thus
$2\Riesz_{2T}^{(1)}-\Riesz_T^{(1)}$ approximates the normalized Stokes
jump $\Ssum$ with exponentially small error. This works even when
$\beta(\alpha)=\infty$. It does not require finding the close pairs
before defining the weighted finite sum.

\subsection{Established ingredients and the novelty boundary}
Continued-fraction approximation and the exponential approximation
exponent are classical. In particular, Sondow's irrationality base has
logarithm $\limsup\log q_{k+1}/q_k$ \cite{sondow}; our $b$ is that
\emph{logarithm}, not the base itself. Divided differences and their
Genocchi--Hermite representation are classical \cite{deboor}. Riesz means
of general Dirichlet series have a substantial theory; the work of
Defant and Schoolmann gives a modern primary reference
\cite{defant-schoolmann}. Richardson extrapolation is ordinary polynomial
interpolation at zero in the variable $1/T$; we prove the needed finite
identities instead of attributing novelty to that procedure.

The additions relative to the inspected repository report are the exact
critical-line reduction, its ordering-sensitive realization theory, the
raw derivative hierarchy, and the sharp cutoff-regularization theorem
with its finite bias and exponential acceleration. These are
model-specific results, not a general solution to all summability
questions for transseries. Their proofs stand independently of a claim
to being the first occurrence of every refinement in the literature.
The computational checks are supplementary evidence, not substitutes
for the proofs.

\section{The Borel model, actions, and block summation}
\label{sec:model}

\subsection{Definitions and conventions}
Fix integers $d\ge2$ and $D\ge d$. We use the slight extension
\begin{equation}
 \widehat H_{\boldsymbol\omega,D}(t)
  =\frac{t^D}{\prod_{j=1}^d\sin(\pi\omega_jt)},\qquad
 \Hh=\{z\in\C:\re z>0\}.
 \label{eq:kernel}
\end{equation}
The singularity at zero is removable. The original model is $D=d$;
allowing $D$ to vary makes differentiation transparent. To make the
transseries connection explicit, expand the analytic germ at zero as
$\widehat H(t)=\sum_{n\ge0}h_nt^n$. In the Borel convention used here it
corresponds to the formal inverse-power series
\begin{equation}
 \widetilde H(z)=\sum_{n\ge0}n!h_nz^{-n-1}.
 \label{eq:formal-borel}
\end{equation}
Its lateral Laplace realizations differ by the exponential residue
corrections studied below. A finite union of positive action lattices
is locally finite and lies in the additive monoid generated by their
positive spacings. Thus the obstruction is not an invalid formal action
support: it is failure of an infinite raw numerical sum at a fixed
argument. Finite-block or Riesz realization repairs the latter without
changing the input actions.

Until the confluence discussion, assume that $\omega_i/\omega_j$ is irrational
for $i\ne j$. All positive poles are then simple. At
$\lambda=n/\omega_j$ their residues are
\begin{equation}
 c_{n/\omega_j}
 =\frac{(-1)^n(n/\omega_j)^D}
 {\pi\omega_j\prod_{\ell\ne j}\sin(\pi\omega_\ell n/\omega_j)}.
 \label{eq:residue-general}
\end{equation}
The action set $\Lambda$ is the union of these $d$ positive lattices.
It is locally finite. Define the increasing-action cutoff by
\begin{equation}
 P_T(z)=\sum_{\substack{\lambda\in\Lambda\\\lambda<T}}
                    c_\lambda\e^{-\lambda z}.
 \label{eq:raw-cutoff}
\end{equation}
Ordinary increasing-action convergence means that $P_T(z)$ has a limit
as $T\to\infty$, with no restriction that $T$ avoid a small gap in a
particular way. The strict inequality fixes values at poles but does
not affect the convergence question.

For two lattices we write the two separate series as
\begin{align}
 r_n(z)&=\frac{(-1)^n n^D\e^{-nz}}{\pi\sin(\pi\alpha n)},
 &R(z)&=\sum_{n\ge1}r_n(z),\label{eq:two-r}\\
 s_m(z)&=\frac{(-1)^m(m/\alpha)^D\e^{-(m/\alpha)z}}
                  {\pi\alpha\sin(\pi m/\alpha)},
 &S(z)&=\sum_{m\ge1}s_m(z).
 \label{eq:two-s}
\end{align}
The capital letters $R,S$ in this display refer only to the separate
series. The calligraphic $\Ssum$ below always denotes the block sum.

\subsection{Uniform blocks: a recalled mechanism with proof}
Fix a compact positive period box
\begin{equation}
 0<m_0\le\omega_j\le M_0,\qquad
 h=\frac{1}{8(d+1)M_0}.
 \label{eq:box}
\end{equation}
Join successive pole occurrences when their gap is less than $h$.
The connected components are called blocks. When a repeated pole is
allowed, its occurrences belong to one block and are counted with
multiplicity.

\begin{lemma}[Geometry and pole-cancelled factors]
\label{lem:block-geometry}
Every block $C$ has at most $d$ occurrences and interval length at most
$(d-1)h$. Consecutive blocks are separated by at least $h$. Their minima
have a uniformly bounded count in every unit interval. Put $a_C=\min C$
and
\begin{equation}
 g_C(t,z)=\e^{-zt}\widehat H_{\boldsymbol\omega,D}(t)
                         \prod_{\lambda\in C}(t-\lambda).
 \label{eq:regular-factor}
\end{equation}
There is a fixed tube around each real block interval on which $g_C$ is
holomorphic. For every compact $K\Subset\Hh$ and every fixed integer
$J\ge0$, constants independent of internal pole gaps satisfy
\begin{equation}
 \sup_{\substack{z\in K\\t\in[\min C,\max C]}}
       |\partial_t^j g_C(t,z)|
 \le C_{K,J}(1+a_C)^D\e^{-\sigma a_C},
 \quad 0\le j\le J,\qquad \sigma=\min_K\re z.
 \label{eq:regular-bound}
\end{equation}
\end{lemma}
\begin{proof}
If $d+1$ successive occurrences belonged to a block, their first $d$
gaps would give a span less than $dh<1/M_0$. Two occurrences would come
from one lattice, contradicting its spacing $1/\omega_j\ge1/M_0$.
This also shows that a block contains at most one occurrence from each
lattice. The separation statement is the definition. A unit interval
contains at most $d(M_0+1)$ pole occurrences, rounded upward if necessary,
which bounds its number of block minima.

Take a tube of radius, for example, $h/4$ around the block interval.
An uncancelled positive zero is at distance at least $h$ from the
interval. The tube also stays away from zero and the negative poles.
For a lattice represented by a node $\lambda=n/\omega_j$, its cancelled
factor is, up to sign,
\[
 \frac{t-\lambda}{\sin(\pi\omega_j(t-\lambda))}.
\]
It is uniformly bounded on a slightly larger tube: the argument remains
strictly inside the distance to the next zero, and its removable value
is $1/(\pi\omega_j)$. For a lattice not represented in the block, the
sine factor is bounded away from zero on that tube. These assertions
follow either from $|\sin(u+\ii v)|^2=\sin^2u+\sinh^2v$ or from
compactness after reduction modulo its period. The period box supplies
uniform constants.

The numerator is bounded by a constant times $(1+a_C)^D$. On the tube,
$|\e^{-zt}|\le C_K\e^{-\sigma a_C}$, since its width and the block
length are fixed and $z$ ranges over $K$. Cauchy's estimates on a nested
tube give every fixed derivative bound in \eqref{eq:regular-bound}.
No step divides by an internal gap.
\end{proof}

For distinct real nodes $\lambda_1,\ldots,\lambda_r$, write
\begin{equation}
 [\lambda_1,\ldots,\lambda_r]f
   =\sum_{i=1}^r\frac{f(\lambda_i)}
           {\prod_{j\ne i}(\lambda_i-\lambda_j)}.
 \label{eq:dd}
\end{equation}
The divided difference has its usual confluent meaning at repeated
nodes. The Genocchi--Hermite formula, with simplex volume $1/(r-1)!$,
gives
\begin{equation}
 |[\lambda_1,\ldots,\lambda_r]f|
  \le\frac{\sup_I|f^{(r-1)}|}{(r-1)!},
 \qquad I=[\min\lambda_i,\max\lambda_i].
 \label{eq:dd-bound}
\end{equation}
This follows by iterating the fundamental theorem of calculus; see also
\cite[formula (52)]{deboor}. The same estimate holds for distinct nodes
when $f\in W^{r-1,\infty}(I)$, with the essential supremum on the right.
Indeed, extend $f$ across the endpoints using its derivative traces,
mollify, apply the smooth estimate, and pass to the uniform limit at
the finitely many nodes. For $r=1$ it is simply a bound on a value.
We will need this Sobolev version at a nonsmooth cutoff, but will not
use it to assign an undefined derivative to a fully confluent node.

\begin{proposition}[Block sum and the Stokes jump]
\label{prop:block-sum}
For a block $C$, let
\begin{equation}
 B_C(z)=\sum_{\lambda\in C\ \mathrm{distinct}}
      \Res_{t=\lambda}\bigl(\e^{-zt}\widehat H(t)\bigr)
       =[\lambda_1,\ldots,\lambda_r]g_C(\cdot,z).
 \label{eq:block-residue}
\end{equation}
Then $\Ssum(z)=\sum_C B_C(z)$ converges absolutely and locally uniformly
on $\Hh$, together with each fixed $z$-derivative. It is independent of
the sufficiently small admissible block threshold. With lower and upper
lateral Laplace integrals $F_-,F_+$,
\begin{equation}
 F_-(z)-F_+(z)=2\pi\ii\Ssum(z).
 \label{eq:stokes}
\end{equation}
All block estimates are uniform on the period box, including confluent
blocks.
\end{proposition}
\begin{proof}
The residue identity follows from the contour formula for divided
differences. Equations \eqref{eq:regular-bound} and \eqref{eq:dd-bound}
majorize $|B_C(z)|$ by $C_K(1+a_C)^D\e^{-\sigma a_C}$. The bound on the
number of minima in a unit interval makes the sum finite. Multiplying
the numerator by any fixed power of $t$ proves the derivative assertion.

For completeness, choose small rays above and below the positive axis
such that $\re(z\e^{\ii\theta})>0$ along the intervening angles.
The Laplace integrals converge at zero because $D\ge d$, and at infinity
because the sine factors grow in either open half-plane. Close the rays
at radii in the middle of the gaps between blocks. These radii stay a
fixed distance from real poles. On the closing arcs, the sine factors
are bounded away from zero near the real axis and grow away from it;
the numerator is at most polynomial. The exponentially decaying
Laplace factor makes the outer arc tend to zero. The residue theorem
gives \eqref{eq:stokes}, with the lower ray traversed outward. Every
admissible blocking has the same two lateral integrals, which proves
independence. This is the predecessor's block argument with $D$ in
place of $d$ \cite{repo-block}.
\end{proof}

\section{Continued fractions and exponentially weighted small divisors}
\label{sec:cf}

\subsection{The approximation facts, including their proofs}
Write $\alpha=[a_0;a_1,a_2,\ldots]$ and use
\begin{align*}
 p_{-2}=0,\quad p_{-1}=1,&\qquad p_k=a_kp_{k-1}+p_{k-2},\\
 q_{-2}=1,\quad q_{-1}=0,&\qquad q_k=a_kq_{k-1}+q_{k-2}.
\end{align*}
The index $k$ always retains this convention; it determines a sign in
the main theorem. Finitely many initial exceptional indices, including
a repeated initial denominator or a zero numerator, may be deleted
from summability statements without renumbering the remaining indices.
Set $\delta_k=\alpha q_k-p_k$.

\begin{lemma}[Continued-fraction estimates]
\label{lem:cf}
For all sufficiently large $k$,
\begin{equation}
 \operatorname{sgn}\delta_k=(-1)^k,\qquad
 q_{k+1}<|\delta_k|^{-1}<q_{k+1}+q_k.
 \label{eq:cf-error}
\end{equation}
For $q_k\le n<q_{k+1}$,
\begin{equation}
 \|n\alpha\|\ge|\delta_k|,
 \qquad \|u\|:=\dist(u,\Z).
 \label{eq:best-approx}
\end{equation}
If integers $p,n$, $n>0$, satisfy $|n\alpha-p|<1/(2n)$, then the reduced
fraction $p/n$ is a convergent. Finally, $q_k$ grows at least
geometrically in $k$ after grouping adjacent indices.
\end{lemma}
\begin{proof}
Let $\tau_{k+1}=a_{k+1}+\theta_k$ be the complete quotient, where
$0<\theta_k<1$. The recurrence gives
\[
 \alpha=\frac{p_k\tau_{k+1}+p_{k-1}}
               {q_k\tau_{k+1}+q_{k-1}},\qquad
 \delta_k=\frac{(-1)^k}{q_k\tau_{k+1}+q_{k-1}},
\]
since $p_kq_{k-1}-p_{k-1}q_k=(-1)^{k-1}$. This proves
\eqref{eq:cf-error} and
$\delta_{k-1}=-\tau_{k+1}\delta_k$.

To prove the best-approximation assertion, suppose $0<n<q_{k+1}$ and
$|n\alpha-p|<|\delta_k|$. The determinant identity writes
$(p,n)=u(p_k,q_k)+v(p_{k-1},q_{k-1})$ with $u,v\in\Z$. Thus
$|u-v\tau_{k+1}|<1$. If $v>0$, this implies $u\ge va_{k+1}$ and
$n\ge vq_{k+1}$, a contradiction. If $v<0$, it implies
$u\le va_{k+1}$ and $n\le vq_{k+1}<0$. If $v=0$, it forces $u=0$.
All possibilities are impossible. Taking the nearest $p$ proves
\eqref{eq:best-approx}; eventually $p_k$ itself is nearest to
$q_k\alpha$.

For the Legendre assertion, first reduce $p/n=P/Q$. Its error is less
than $1/(2Q^2)$. Choose $k$ with $q_k\le Q<q_{k+1}$. Unless
$P/Q=p_k/q_k$, best approximation and the determinant bound give
\[
 1\le|Pq_k-Qp_k|
 \le q_k|Q\alpha-P|+Q|\delta_k|
 \le2Q|Q\alpha-P|<1,
\]
a contradiction. Initial small-denominator cases follow from the same
continued-fraction identities, or can be omitted in our tail arguments.
Finally $q_{k+2}\ge q_{k+1}+q_k\ge2q_k$ once denominators increase.
\end{proof}

\begin{corollary}[Exponential exponent and reciprocity]
\label{cor:beta}
For every positive irrational $\alpha$,
\begin{equation}
 \beta(\alpha)=\limsup_k\frac{\log q_{k+1}}{q_k},
 \qquad \alpha\beta(1/\alpha)=\beta(\alpha),
 \label{eq:beta-cf}
\end{equation}
including the values zero and infinity.
\end{corollary}
\begin{proof}
Evaluating at $n=q_k$ and using \eqref{eq:cf-error} proves one inequality
in the first formula. For $q_k\le n<q_{k+1}$, best approximation bounds
$n^{-1}\log(1/\|n\alpha\|)$ above by
$q_k^{-1}\log(q_{k+1}+q_k)$. The difference between this and
$q_k^{-1}\log q_{k+1}$ tends to zero. This proves equality. It agrees
with the classical irrationality-base formula \cite{sondow}.

Apart from initial terms, the convergents of $1/\alpha$ are $q_k/p_k$.
Since $p_k/q_k\to\alpha$ and
$\log p_{k+1}=\log q_{k+1}+O(1)$, the second equality follows, with
extended-real limits understood.
\end{proof}

\subsection{An exact positive-series and term-test reduction}

\begin{lemma}[Weighted denominator test]
\label{lem:weighted-test}
Let $x>0$ and $D\ge0$ be fixed, and set
$A_k(x)=q_k^Dq_{k+1}\e^{-xq_k}$. Then
\begin{align}
 \sum_{n\ge1}\frac{n^D\e^{-xn}}{|\sin(\pi n\alpha)|}<\infty
 &\quad\Longleftrightarrow\quad \sum_k A_k(x)<\infty,
 \label{eq:positive-test}\\
 \frac{n^D\e^{-xn}}{|\sin(\pi n\alpha)|}\longrightarrow0
 &\quad\Longleftrightarrow\quad A_k(x)\longrightarrow0.
 \label{eq:term-test}
\end{align}
\end{lemma}
\begin{proof}
Use $2\|u\|\le|\sin(\pi u)|\le\pi\|u\|$. For the integer block
$q_k\le n<q_{k+1}$, Lemma~\ref{lem:cf} gives the upper bound
\begin{align*}
 \sum_{n=q_k}^{q_{k+1}-1}
       \frac{n^D\e^{-xn}}{|\sin(\pi n\alpha)|}
 &\le\frac{q_{k+1}+q_k}{2}
                \sum_{n\ge q_k}n^D\e^{-xn}\\
 &\le C_{x,D}q_k^Dq_{k+1}\e^{-xq_k}.
\end{align*}
Here $n=q_k+j$ and $(q_k+j)^D\le q_k^D(1+j)^D$ bound the last sum.
The single convergent index $n=q_k$ is bounded below by
$A_k(x)/\pi$. The denominator blocks partition a tail of the positive
integers. These inequalities prove \eqref{eq:positive-test}. They also
bound every term in the $k$th block by $C_{x,D}A_k(x)$, proving both
directions of \eqref{eq:term-test}.
\end{proof}

This lemma already identifies the absolute convergence criterion on a
critical line. Conditional convergence requires more information: the
near-cancelling residues must be matched and their signs retained. That
is the purpose of the next section.

\section{Critical-line extraction and the complete ordering criterion}
\label{sec:critical}
Assume throughout this section that $D\ge2$, $\alpha>0$ is irrational,
and $0<b=\beta(\alpha)<\infty$. Define
\begin{equation}
 A_k=q_k^Dq_{k+1}\e^{-bq_k},\qquad
 \eta_k=(-1)^{p_k+q_k+k},\qquad
 \lambda_k'=p_k/\alpha=q_k-\delta_k/\alpha.
 \label{eq:critical-data}
\end{equation}
The two actions $q_k$ and $\lambda_k'$ will be called the primitive
convergent pair.

\subsection{The absolutely summable remainder}

\begin{lemma}[Nonprimitive residues are summable]
\label{lem:nonprimitive}
After deleting $r_{q_k}$ from the first lattice and $s_{p_k}$ from the
second lattice, the remaining residue terms on $\re z=b$ are absolutely
summable. The sums of their absolute values are independent of $\im z$.
\end{lemma}
\begin{proof}
First consider an integer $n$ in the first lattice. If
$\|n\alpha\|\ge1/(2n)$, then
$|r_n(b+\ii y)|\le n^{D+1}\e^{-bn}/\pi$, and these bounds sum.
Otherwise choose the nearest integer $p$ and reduce $p/n$. By
Lemma~\ref{lem:cf},
\[
 n=\ell q_k,\qquad p=\ell p_k,\qquad
 \|n\alpha\|=\ell|\delta_k|
\]
for a unique reduced convergent and a positive integer $\ell$.
A nonprimitive term has $\ell\ge2$. Its absolute value is at most
\[
 C q_k^D(q_{k+1}+q_k)\ell^{D-1}\e^{-b\ell q_k}.
\]
The sum over $\ell\ge2$ is at most
$C_{b,D}q_k^D(q_{k+1}+q_k)\e^{-2bq_k}$.
Because $\limsup\log q_{k+1}/q_k=b$, eventually
$q_{k+1}\le\e^{3bq_k/2}$. Thus the last bound is summable in $k$,
using geometric growth of $q_k$. Finitely many early $k$ cause no
problem, since their sums over $\ell$ converge.

Apply the same proof to $1/\alpha$ with the Laplace parameter
$b/\alpha=\beta(1/\alpha)$. Its primitive denominators are $p_k$, apart
from initial terms. The constant factor $\alpha^{-D-1}$ in the second
lattice does not affect summability. Absolute values remove $y$ from
both arguments.
\end{proof}

\begin{theorem}[Critical-line extraction modulo \(\ell^1\)]
\label{thm:extraction}
For every $Y<\infty$ there are functions $u_k,v_k$ on $[-Y,Y]$ such that
\begin{align}
 r_{q_k}(b+\ii y)
    &=\frac{\eta_k}{\pi^2}A_k\e^{-\ii yq_k}+u_k(y),
       \label{eq:r-extraction}\\
 s_{p_k}(b+\ii y)
    &=-\frac{\eta_k}{\pi^2}A_k\e^{-\ii yq_k}+v_k(y),
       \label{eq:s-extraction}\\
 \sum_k\sup_{|y|\le Y}\bigl(|u_k(y)|+|v_k(y)|\bigr)&<\infty.
       \label{eq:ellone}
\end{align}
Together with Lemma~\ref{lem:nonprimitive}, these formulas account for
all but finitely many raw residues.
\end{theorem}
\begin{proof}
Let $F_z(t)=t^D\e^{-zt}$. The first residue has the form
\[
 r_{q_k}(z)=\frac{(-1)^{p_k+q_k}F_z(q_k)}{\pi\sin(\pi\delta_k)}.
\]
For small real $\delta$, $1/\sin(\pi\delta)=1/(\pi\delta)+O(|\delta|)$.
Moreover
$|\delta_k|^{-1}=q_{k+1}+\theta_kq_k$, $0<\theta_k<1$.
Using $\operatorname{sgn}\delta_k=(-1)^k$ gives
\eqref{eq:r-extraction}, with
\begin{equation}
 \sup_{|y|\le Y}|u_k(y)|\le C_Yq_k^{D+1}\e^{-bq_k}
 \label{eq:error-majorant}
\end{equation}
for large $k$.

For the companion residue, an exact sine identity gives
\[
 s_{p_k}(z)
 =-\frac{(-1)^{p_k+q_k}F_z(q_k-\delta_k/\alpha)}
             {\pi\alpha\sin(\pi\delta_k/\alpha)}.
\]
Replace $1/(\pi\alpha\sin(\pi\delta_k/\alpha))$ by
$1/(\pi^2\delta_k)$; the error is $O(|\delta_k|)$.
The mean-value integral on the short real segment yields, uniformly for
$|y|\le Y$,
\[
 |F_{b+\ii y}(q_k-\delta_k/\alpha)-F_{b+\ii y}(q_k)|
 \le C_Y|\delta_k|q_k^D\e^{-bq_k}.
\]
After division by $\delta_k$, this has a summable majorant. Replacing
$|\delta_k|^{-1}$ by $q_{k+1}$ adds at most the bound in
\eqref{eq:error-majorant}. This proves \eqref{eq:s-extraction} and
\eqref{eq:ellone}.
\end{proof}

\subsection{The classification theorem}

\begin{theorem}[All ordinary critical-line convergence types]
\label{thm:classification}
At every $z=b+\ii y$, the following statements hold.
\begin{enumerate}[label=\textup{(\roman*)}]
\item The increasing-action series $P_T(z)$ converges if and only if
$A_k\to0$. Its value, when it converges, is $\Ssum(z)$.
\item It converges absolutely, equivalently unconditionally as a complex
series, if and only if $\sum_k A_k<\infty$.
\item The two separately indexed lattice series in
\eqref{eq:two-r}--\eqref{eq:two-s} both converge if and only if
\begin{equation}
 \sum_k\eta_k A_k\e^{-\ii yq_k}
 \quad\hbox{converges.}
 \label{eq:separate-criterion}
\end{equation}
Either separate lattice series converges precisely under this same
criterion. When they converge, their sum equals $\Ssum(z)$.
\end{enumerate}
If $A_k\to0$, increasing-action convergence is uniform on each compact
segment of $\re z=b$, and on each compact rectangle
$b\le\re z\le B$, $|\im z|\le Y$.
\end{theorem}
\begin{proof}
Use the two-lattice blocks of Section~\ref{sec:model}; each has at most
two nodes and its nodes are consecutive in action order. At a cutoff
$T$, all retained blocks are complete except possibly the last one.
The sum of the complete blocks tends locally uniformly to $\Ssum$.
The incomplete block, if present, contributes exactly one raw residue.
Therefore $P_T(z)$ converges to $\Ssum(z)$ whenever all individual
residues tend to zero. Conversely, convergence of an ordinary series
forces its individual terms to tend to zero.

Lemma~\ref{lem:weighted-test} identifies the first-lattice term test
with $A_k\to0$. For the second lattice, the same conclusion follows
from Lemma~\ref{lem:nonprimitive} and Theorem~\ref{thm:extraction}.
This proves (i). All term moduli on the critical line are independent
of $y$; the same bounds decrease on moving to the right. The block
sum is locally uniform on $\Hh$. These observations prove the stated
uniformity.

The absolute convergence equivalence follows either from
Lemma~\ref{lem:weighted-test} and the companion estimates, or directly
from extraction: an $\ell^1$ perturbation cannot change whether the
positive sum of the leading moduli $A_k/\pi^2$ is finite. This proves
(ii). Finally the deleted remainders and $u_k,v_k$ are absolutely
summable, whereas the retained leading subseries are opposite scalar
multiples of \eqref{eq:separate-criterion}. This proves (iii).
If the separate sums converge, their increasing-action interleaving
has their sum as its limit, which (i) identifies with $\Ssum$.
\end{proof}

\begin{remark}[What the separate criterion does and does not say]
Equation \eqref{eq:separate-criterion} is a necessary-and-sufficient
criterion in explicit continued-fraction data, not a sufficient
Diophantine bound. It reduces the full lattice sums to one sparse
signed exponential series, with an absolutely summable remainder.
For general $y$ it retains a genuine phase-convergence question. The
increasing-action criterion has no such residual phase dependence.
A detailed classification of the set of imaginary heights satisfying
\eqref{eq:separate-criterion} is a further problem, not hidden inside
an assertion of absolute convergence.
\end{remark}

\subsection{The zero and infinite endpoints}
The positive finite case is the only nontrivial ordinary boundary
case. If $\beta(\alpha)=0$ and $\re z=0$, then
\[
 |r_n(z)|=\frac{n^D}{\pi|\sin(\pi n\alpha)|}\ge\frac{n^D}{\pi}.
\]
Thus both increasing-action and separate ordinary interpretations fail
the term test on that entire boundary line. If $\beta(\alpha)=\infty$,
there is no finite critical boundary, and the raw terms fail to vanish
at every fixed $\re z>0$. The open-half-plane statements also follow
from \eqref{eq:beta-cf}: the exponential root test gives absolute
convergence for $\re z>\beta(\alpha)$ and nonvanishing terms below it.
Together these facts and Theorem~\ref{thm:classification} complete the
ordinary convergence classification requested in the predecessor's
Question~1. No assertion about analytic continuation of $\Ssum$ across
$\re z=0$ is needed.

\section{Prescribing the boundary behavior}
\label{sec:realization}

\subsection{A continued-fraction amplitude construction}

\begin{theorem}[Prescribed critical amplitudes]
\label{thm:construction}
Fix $b>0$ and an integer $D\ge2$. Start with a finite continued-fraction
prefix. At stage $k$, choose a positive target $t_k$, and choose the next
partial quotient within a bounded distance of
\begin{equation}
 X_k=\frac{t_k\e^{bq_k}}{q_k^{D+1}},
 \label{eq:target-quotient}
\end{equation}
subject to positivity. One may require a specified parity, since the
distance to a positive integer of that parity is bounded once $X_k$ is
large. Suppose the chosen target prescription satisfies
\begin{equation}
 |\log t_k|=o(q_k)
 \label{eq:target-condition}
\end{equation}
along the resulting construction. Then the resulting infinite continued
fraction is irrational, has $\beta(\alpha)=b$, and satisfies
\begin{equation}
 A_k=t_k+O(q_k^{D+1}\e^{-bq_k}),\qquad
 A_k/t_k\longrightarrow1,\qquad
 \sum_k|A_k-t_k|<\infty.
 \label{eq:target-error}
\end{equation}
The limits are unaffected by any finite initial choices.
\end{theorem}
\begin{proof}
The recurrence and bounded rounding error give
\[
 q_{k+1}=\frac{t_k\e^{bq_k}}{q_k^D}+O(q_k),
\]
because $q_{k-1}\le q_k$ eventually. Multiplying by
$q_k^D\e^{-bq_k}$ gives the first part of \eqref{eq:target-error}.
Condition \eqref{eq:target-condition} ensures $X_k\to\infty$, so the
rounding rule is feasible eventually and its error is negligible
relative to $t_k$. Geometric denominator growth makes the error
absolutely summable. Taking logarithms yields
\[
 \frac{\log q_{k+1}}{q_k}
   =b+\frac{\log t_k-D\log q_k+o(1)}{q_k}\longrightarrow b.
\]
Corollary~\ref{cor:beta} proves the exponent assertion. An infinite
simple continued fraction with positive partial quotients is irrational:
its consecutive reduced convergents are distinct, with errors tending
to zero and alternating signs. Equivalently, the complete-quotient
formulas rule out a rational limit with a fixed denominator.
\end{proof}

The prescription can be an explicit function of the already known pair
$(k,q_k)$; hence the construction is inductive, not an equation for an
unknown irrational. For example $t_k=k^{-r}$ and
$t_k=q_k^{-r}/k^s$ satisfy \eqref{eq:target-condition}. For the concrete
examples, choose $a_0=1$, $a_1\ge2$, so $1<\alpha<2$.

\begin{corollary}[Every ordinary boundary type at every positive abscissa]
\label{cor:types}
At each prescribed $b>0$ there are irrational period ratios exhibiting
all three cases in Theorem~\ref{thm:classification}. Targets $k^{-2}$,
$k^{-1}$, and $1$ give, respectively, absolute convergence, conditional
increasing-action convergence, and failure of the term test. The target
$k$ produces an unbounded subsequence of individual residue moduli.
\end{corollary}
\begin{proof}
The first two claims follow from \eqref{eq:target-error}; the summable
rounding errors cannot change the relevant positive-series tests.
For the other two, $A_k\sim t_k$ and the extraction theorem give the
asserted nonvanishing or unbounded primitive residues.
\end{proof}

\subsection{An increasing-action sum with divergent separate sums}

\begin{example}[Opposite divergent lattice sums]
\label{ex:opposite}
Set $a_0=1$, choose all $a_j$ for $j\ge1$ even, and use the even
rounding version of Theorem~\ref{thm:construction} with $t_k=1/k$ in
the tail. Modulo two, the convergent vectors alternate as
\[
 (p_k,q_k)\equiv
 \begin{cases}(1,1),&k\text{ even},\\(1,0),&k\text{ odd}.
 \end{cases}
\]
This follows directly from
$(p_{k+1},q_{k+1})\equiv(p_{k-1},q_{k-1})\pmod2$.
Consequently $\eta_k=1$ for every $k$. At $z=b$, the first separate
lattice series equals a positive harmonic-type divergent series plus
a convergent remainder and tends to $+\infty$; the second tends to
$-\infty$. Nevertheless, $A_k\to0$ and $\sum A_k=\infty$, so their
increasing-action interleaving converges conditionally to $\Ssum(b)$.
\end{example}

This example is not an assignment of a finite value to
$+\infty-\infty$. The increasing-action series is a separately defined
ordinary series with its own partial sums. The ordering makes its
close-pair cancellations occur before the next large-scale contribution,
and the remaining unfinished term tends to zero.

\begin{example}[Both separate lattice sums conditionally converge]
\label{ex:separate-conditional}
Choose $a_0=1$, $a_1$ even, $a_2$ odd, and all later partial quotients
even. For $k\ge1$, the parity vectors alternate between $(1,0)$ and
$(0,1)$. Therefore $\eta_k=(-1)^{k+1}$. Choose the tail quotients by
Theorem~\ref{thm:construction} with $t_k=1/k$, respecting even parity.
Then
\[
 \sum_k\eta_k A_k
   =\sum_k\frac{(-1)^{k+1}}{k}
        +\text{an absolutely convergent series}
\]
converges. Both separate lattice series converge at $b$, neither
absolutely. Their increasing-action interleaving also converges.
The argument uses a summable perturbation of the alternating harmonic
series, not an unproved monotonicity assertion about $A_k$.
\end{example}

\subsection{A two-point cluster set for a divergent raw expansion}

\begin{proposition}[Bounded spikes of prescribed height]
\label{prop:spikes}
Fix $L>0$. Use the parity pattern of
Example~\ref{ex:separate-conditional} and the constant target $t_k=L$.
At the real critical point $z=b$, the increasing-action partial sums
have exactly two finite limit points:
\begin{equation}
 \Ssum(b)\quad\text{and}\quad\Ssum(b)+\frac{L}{\pi^2}.
 \label{eq:cluster-set}
\end{equation}
They are bounded but do not converge. With the same parity pattern and
$t_k=k$, a subsequence of partial sums tends to $+\infty$, while
block-complete partial sums still tend to $\Ssum(b)$.
\end{proposition}
\begin{proof}
For large $k$, $q_k$ and $\lambda_k'$ form one two-node block, since
$|\delta_k|\to0$. If $k$ is even, $\delta_k>0$ and the companion
$\lambda_k'<q_k$ is first. If $k$ is odd, the integer action $q_k$ is
first. In either case the first residue has sign
$(-1)^{p_k+q_k+1}=+1$ in its leading term and equals
$A_k/\pi^2+o(1)$. The second residue cancels that leading term.

Every nonprimitive residue belongs to an absolutely summable family
by Lemma~\ref{lem:nonprimitive}, and hence tends to zero individually.
Every primitive residue belongs to its indicated pair, apart from
finitely many indices. A raw cutoff differs from its last completed
block sum by either zero, a vanishing nonprimitive term, or the first
term of a primitive pair. Completed block sums tend to $\Ssum(b)$.
For $t_k=L$, the only possible limit points are the two in
\eqref{eq:cluster-set}, and cuts after and inside the primitive pairs
realize both. For $t_k=k$, the same cuts inside the pairs diverge to
$+\infty$. No assertion is being made here about the cluster sets at
nonzero imaginary heights.
\end{proof}

\section{The derivative hierarchy of a raw transseries}
\label{sec:derivatives}
A normally convergent block sum is holomorphic, but this does not
license differentiating a different, conditionally convergent raw
representation term by term at a boundary point.

\begin{theorem}[Exact threshold for each formal derivative]
\label{thm:derivatives}
Under the hypotheses of Section~\ref{sec:critical}, fix an integer
$j\ge0$. The increasing-action series formed by differentiating every
raw term $j$ times converges at $b+\ii y$ if and only if
\begin{equation}
 q_k^jA_k\longrightarrow0.
 \label{eq:derivative-test}
\end{equation}
It converges absolutely if and only if $\sum_k q_k^jA_k<\infty$.
Whenever it converges, its value is $\Ssum^{(j)}(b+\ii y)$, uniformly
on compact segments of the critical line. The separate derivative
series converge precisely when
\begin{equation}
 \sum_k\eta_k q_k^jA_k\e^{-\ii yq_k}
 \quad\hbox{converges.}
 \label{eq:separate-derivative}
\end{equation}
\end{theorem}
\begin{proof}
Differentiating $c_\lambda\e^{-\lambda z}$ multiplies it by
$(-\lambda)^j$. Apart from the common factor $(-1)^j$, the result is
exactly the degree-$D+j$ residue family. Its critical amplitude is
$q_k^jA_k$. Apply the extraction and classification theorems with
$D+j$ in place of $D$. Proposition~\ref{prop:block-sum} identifies
the corresponding block sum with the derivative of $\Ssum$.
\end{proof}

\begin{corollary}[Arbitrarily prescribed finite differentiation depth]
\label{cor:derivative-examples}
Fix an integer $r\ge0$. There are period ratios with any prescribed
$b>0$ for which derivatives of orders below $r$ have absolutely
convergent raw series, order $r$ has a conditionally convergent
increasing-action raw series, and order $r+1$ fails the term test.
There are also examples with absolute raw convergence at every
derivative order on $\re z=b$, while the absolute convergence abscissa
remains exactly $b$.
\end{corollary}
\begin{proof}
For the first statement prescribe
$t_k=q_k^{-r}/k$. Then $q_k^jA_k$ differs from $q_k^{j-r}/k$ by a
summable, exponentially small sequence for every fixed $j$.
For $j<r$, geometric growth of $q_k$ gives summability. At $j=r$
the amplitudes tend to zero but have divergent sum. At $j=r+1$,
$q_k/k\to\infty$.

For the second statement prescribe $t_k=\e^{-\sqrt{q_k}}$.
Condition \eqref{eq:target-condition} holds, so $\beta(\alpha)=b$.
For each fixed $j$, $\sum q_k^j\e^{-\sqrt{q_k}}<\infty$.
Theorem~\ref{thm:derivatives} gives the assertion.
\end{proof}

The failure in the first statement belongs to the representation, not
to the analytic function: $\Ssum$ has derivatives of all orders near
every point with positive real part. At a high enough order the raw
termwise derivative ceases to be an ordinary convergent series, but the
block derivative remains valid. This separates analytic regularity from
raw-sector regularity without changing the formal action set.

\section{Riesz cutoff regularization without Diophantine assumptions}
\label{sec:riesz}
We now return to $d$ lattices, the period box \eqref{eq:box}, and
$D\ge d$. No approximation exponent is assumed finite. Positive poles
are simple in the principal theorem; rational confluence is treated
separately in Section~\ref{sec:sharpness}.

\subsection{A weighted cutoff that can split a block}
For an integer $m\ge1$, define the finite Riesz mean
\begin{equation}
 \Riesz_T^{(m)}(z)=\sum_{\lambda<T}
        \left(1-\frac{\lambda}{T}\right)^m c_\lambda\e^{-\lambda z},
 \qquad T>0.
 \label{eq:riesz-mean}
\end{equation}
This is the usual Riesz weight for a Dirichlet series indexed by the
frequency $\lambda$; see \cite{defant-schoolmann} for general context.
Our proof is a direct local argument for the present meromorphic
kernel. In particular, it does not assume that the raw series has a
nonempty half-plane of convergence.

Let $w_{T,m}(t)=T^{-m}(T-t)_+^m$ on the real line, where
$(u)_+=\max(u,0)$. For a simple block $C=\{\lambda_1,\ldots,\lambda_r\}$,
its contribution to \eqref{eq:riesz-mean} is exactly
\begin{equation}
 [\lambda_1,\ldots,\lambda_r]\bigl(g_C(\cdot,z)w_{T,m}\bigr).
 \label{eq:partial-block}
\end{equation}
The identity remains valid when the cutoff lies between nodes of that
block. It is the central reason the pole gaps need not enter the bound.

\begin{lemma}[A cutoff through a cluster]
\label{lem:cutoff-cluster}
Suppose a block of $r\le d$ distinct nodes straddles $T\ge1$ and
$m\ge r-1$. For $K\Subset\Hh$,
\begin{equation}
 \sup_{z\in K}
 \left|[\lambda_1,\ldots,\lambda_r]
             \bigl(g_C(\cdot,z)w_{T,m}\bigr)\right|
 \le C_K T^{D-m}\e^{-\sigma T},\qquad \sigma=\min_K\re z.
 \label{eq:split-bound}
\end{equation}
The same estimate holds with the ordinary polynomial
$T^{-m}(T-t)^m$ in place of $w_{T,m}$ on that block.
\end{lemma}
\begin{proof}
The block interval has bounded length $W=(d-1)h$, so $|t-T|\le W$
there. For $0\le j\le r-1\le m$, the weak derivatives satisfy
\[
 |w_{T,m}^{(j)}(t)|
 \le (m)_jT^{-m}W^{m-j}
\]
almost everywhere, with harmless replacement by a larger constant if
$W<1$. Here $(m)_j=m!/(m-j)!$. In particular,
$w_{T,m}\in W^{r-1,\infty}$, including the endpoint case $m=r-1$.
Leibniz's rule and \eqref{eq:regular-bound} bound the $(r-1)$st weak
derivative of the product by
$C_KT^{-m}(1+T)^D\e^{-\sigma T}$. The Sobolev divided-difference
bound \eqref{eq:dd-bound} proves \eqref{eq:split-bound}.
For the ordinary polynomial the same argument is smooth and identical.
\end{proof}

\subsection{The finite derivative bias and the exponential remainder}

\begin{theorem}[Uniform Riesz recovery and its exact bias]
\label{thm:riesz}
Let $m\ge d-1$ be an integer. For every compact $K\Subset\Hh$, there
is a constant depending only on $K,d,D,m,m_0,M_0$ such that, for all
$T\ge1$ and every period vector in the box with simple positive poles,
\begin{equation}
 \sup_{z\in K}\left|
 \Riesz_T^{(m)}(z)
  -\sum_{j=0}^m\binom{m}{j}\frac{\Ssum^{(j)}(z)}{T^j}
                 \right|
 \le C_KT^{D-m}\e^{-\sigma T},\qquad \sigma=\min_K\re z.
 \label{eq:riesz-jet}
\end{equation}
Consequently $\Riesz_T^{(m)}\to\Ssum$ locally uniformly on $\Hh$.
The bound does not depend on any internal pole separation or on an
exponential approximation exponent.
\end{theorem}
\begin{proof}
First weight every complete block by the ordinary polynomial
$(1-t/T)^m$, even for blocks above the cutoff. Its normally convergent
sum is
\begin{align}
 Q_T(z)
 &=\sum_C[\lambda_1,\ldots,\lambda_r]
             \bigl(g_C(\cdot,z)(1-t/T)^m\bigr)\notag\\
 &=\sum_{j=0}^m\binom{m}{j}\frac{\Ssum^{(j)}(z)}{T^j}.
 \label{eq:full-polynomial}
\end{align}
Indeed the coefficient $(-t)^j$ in the polynomial is precisely the
factor produced by differentiating $\e^{-zt}$ $j$ times. The derivative
normal convergence in Proposition~\ref{prop:block-sum} justifies the
finite expansion and the block summation.

For blocks wholly below $T$, the polynomial weight equals the truncated
weight. There is at most one block straddling $T$. Both of its weighted
contributions are bounded by Lemma~\ref{lem:cutoff-cluster}, since
$m\ge d-1\ge r-1$.

It remains to bound polynomial-weighted blocks wholly above $T$.
For such a block with minimum $a\ge T$, Leibniz's rule, the fixed
interval length, and \eqref{eq:regular-bound} give
\begin{equation}
 \left|[\lambda_1,\ldots,\lambda_r]
       \bigl(g_C(\cdot,z)(1-t/T)^m\bigr)\right|
 \le C_KT^{-m}(1+a)^D(1+a-T)^m\e^{-\sigma a}.
 \label{eq:above-cutoff}
\end{equation}
Bin the minima in $[T+\ell,T+\ell+1)$, using their bounded count.
Because $T\ge1$, the sum of \eqref{eq:above-cutoff} is bounded by
\[
 C_KT^{D-m}\e^{-\sigma T}
       \sum_{\ell\ge0}(1+\ell)^{D+m}\e^{-\sigma\ell}.
\]
The last series is finite. Absorb it in the constant. Combining this
tail with the possible split block proves \eqref{eq:riesz-jet}.
The derivative terms of positive order have factors $T^{-j}$ and are
uniformly bounded on $K$, which gives convergence to $\Ssum$.
\end{proof}

\begin{remark}[What uniformity means]
The period box may contain rational coincidences, while the principal
formula \eqref{eq:riesz-mean} is initially evaluated on its simple-pole
locus. The estimate is uniform as those coincidences are approached.
It asserts a bound, not continuity of every finite cutoff at a
coincidence situated exactly at $T$. That distinction is essential at
minimal order. Also, the power $T^{D-m}$ may be negative when $m>D$;
there is no restriction $m\le D$.
\end{remark}

\subsection{The two-lattice Ces\`aro interpretation}
For simple poles, a finite interchange gives the exact identity
\begin{equation}
 \Riesz_T^{(1)}(z)=\frac1T\int_0^T P_u(z)\dd u.
 \label{eq:cesaro}
\end{equation}
Thus for two lattices the continuous first Ces\`aro mean of the raw
cutoff converges everywhere in $\Hh$, even when the raw cutoff fails
the term test there. A huge unmatched residue persists only over the
very small interval between its two actions. The factor
$(1-\lambda/T)$ integrates that interval rather than estimating its
height alone. Equation \eqref{eq:riesz-jet} is stronger than mere
Ces\`aro convergence: it gives the exact $T^{-1}$ bias and an
exponentially small remaining error.

This interpretation does not say that a naive floating-point evaluation
of \eqref{eq:riesz-mean} is well conditioned. Tiny sine denominators
and cancellation can still require high precision or stable divided
differences. The theorem controls the mathematical finite sum, not the
roundoff error of an arbitrary implementation.

\section{Sharp cutoff order, confluence, and exponential acceleration}
\label{sec:sharpness}

\subsection{The gap-independent threshold is sharp}

\begin{theorem}[Minimal order for uniform collision control]
\label{thm:sharp}
For $d\ge2$, an integer order $m<d-1$ cannot give a bound uniform
across all simple period vectors in a fixed compact positive box, even
at a fixed positive cutoff and a fixed $z\in\Hh$. More precisely, take
\begin{equation}
 \lambda_1=1-\varepsilon,\qquad
 \lambda_j=1+(j-1)\varepsilon\quad(2\le j\le d),\qquad
 \omega_j=\lambda_j^{-1},\qquad T=1,
 \label{eq:collision-family}
\end{equation}
with positive irrational $\varepsilon\to0$. Then the positive pole sets
are pairwise disjoint and
\begin{equation}
 \Riesz_1^{(m)}(z)
 =-\frac{\e^{-z}}{\pi^d\prod_{j=2}^d j}
       \varepsilon^{m-d+1}(1+O_z(\varepsilon)).
 \label{eq:sharp-scaling}
\end{equation}
Thus $d-1$ is the minimal integer order for the gap-independent uniform
control proved in Theorem~\ref{thm:riesz}.
\end{theorem}
\begin{proof}
For small $\varepsilon$, only the first pole $\lambda_1$ is below
$T=1$; even its second lattice multiple is above one. Distinct ratios
$(1+c_i\varepsilon)/(1+c_j\varepsilon)$ with different rational
$c_i,c_j$ cannot be rational when $\varepsilon$ is irrational, so the
positive pole sets are disjoint.

At $\lambda_1$, the derivative of its own sine factor is
$-\pi\omega_1$, while for $j\ge2$,
\[
 \sin(\pi\omega_j\lambda_1)
   =\pi j\varepsilon(1+O(\varepsilon)).
\]
Also $\lambda_1^D=1+O(\varepsilon)$ and
$\e^{-z\lambda_1}=\e^{-z}(1+O_z(\varepsilon))$.
The residue formula, multiplied by
$(1-\lambda_1)^m=\varepsilon^m$, gives
\eqref{eq:sharp-scaling}. Its coefficient is nonzero, proving the
claimed failure for $m<d-1$.
\end{proof}

The sharpness just proved is \emph{parameter-uniform}. It must not be
misread as a proof that every fixed $d$-lattice input needs order $d-1$.
Some inputs already have absolutely convergent raw expansions. For
$d=2$, examples with $\beta=\infty$ do show that order zero can fail
everywhere in $\Hh$ for one fixed irrational input. A sharp fixed-input
classification for more lattices is a separate question.

\subsection{Boundedness is not cutoff continuity}

\begin{proposition}[One additional order for full cutoff confluence]
\label{prop:confluence}
At a block of multiplicity at most $d$, the finite-cutoff expressions
have gap-independent bounds when $m\ge d-1$. If $m\ge d$, they also
extend continuously as the nodes coalesce, including when their common
limit equals $T$. At $m=d-1$, continuity can fail at that cutoff.
\end{proposition}
\begin{proof}
The boundedness is Lemma~\ref{lem:cutoff-cluster}. If $m\ge d$, then
$t\mapsto(T-t)_+^m$ is $C^{d-1}$. Divided differences of order at most
$d-1$ of the smooth regular factor times this weight have their
continuous confluent limits. This follows from the simplex formula
and continuity of the required derivative, even when the limiting
node equals the cutoff.

For $m=d-1$, consider a $d$-node cluster approaching $T$ with a limiting
regular factor $g(T,z)\ne0$. If all nodes approach from above, its
truncated weighted contribution is zero. If all approach from below,
its limit is
\[
 \frac{1}{(d-1)!}\left.
    \partial_t^{d-1}\bigl(g(t,z)(1-t/T)^{d-1}\bigr)
       \right|_{t=T^-}
   =\frac{(-1)^{d-1}}{T^{d-1}}g(T,z),
\]
which is nonzero. Such configurations are obtained by taking the first
poles of the sine lattices close to $T$. Intermediate configurations
can give other limits, as the divided difference of a truncated power
depends on how the cutoff cuts the collapsing nodes.
\end{proof}

At a coincident pole away from the cutoff, the correct definition is
\begin{equation}
 \Res_{t=\lambda}\left((1-t/T)^m\e^{-zt}\widehat H(t)\right),
 \qquad \lambda<T,
 \label{eq:multiple-weight}
\end{equation}
not $(1-\lambda/T)^m$ times the residue of
$\e^{-zt}\widehat H(t)$. The two formulas agree at a simple pole but
need not agree at a multiple pole, because derivatives of the weight
contribute. For $m\ge d$, excluding a pole exactly at $T$ is compatible
with the continuous extension, since all derivatives through order
$d-1$ of the cutoff weight vanish there. At $m=d-1$ no convention at
that exact pole can make every approach continuous. Uniform exponential
error control survives this local ambiguity, but finite-cutoff
continuity does not.

\subsection{Removing the bias with finitely many cutoffs}

\begin{corollary}[Exponentially accurate extrapolated action cuts]
\label{cor:acceleration}
For $m\ge d-1$, define
\begin{align}
 w_\ell&=\frac{(-1)^{m+1-\ell}\ell^m}
                   {(\ell-1)!(m+1-\ell)!},
       &&1\le\ell\le m+1,\label{eq:weights}\\
 \Accel_T^{(m)}(z)&=\sum_{\ell=1}^{m+1}
                          w_\ell\Riesz_{\ell T}^{(m)}(z).
       \label{eq:acceleration}
\end{align}
Then, with the same uniformity as in Theorem~\ref{thm:riesz},
\begin{equation}
 \sup_{z\in K}|\Accel_T^{(m)}(z)-\Ssum(z)|
      \le C_KT^{D-m}\e^{-\sigma T}.
 \label{eq:accelerated-bound}
\end{equation}
In particular,
\begin{align}
 \Accel_T^{(1)}&=2\Riesz_{2T}^{(1)}-\Riesz_T^{(1)},\label{eq:accel-one}\\
 \Accel_T^{(2)}&=\tfrac12\Riesz_T^{(2)}-4\Riesz_{2T}^{(2)}
                         +\tfrac92\Riesz_{3T}^{(2)}.
 \label{eq:accel-two}
\end{align}
\end{corollary}
\begin{proof}
The weights are the Lagrange interpolation weights for evaluating a
polynomial of degree at most $m$ at zero from its values at
$1,1/2,\ldots,1/(m+1)$. Indeed the weight at $1/\ell$ is
\[
 \prod_{r\ne\ell}\frac{-1/r}{1/\ell-1/r}
 =\frac{(-1)^{m+1-\ell}\ell^m}{(\ell-1)!(m+1-\ell)!}.
\]
Therefore $\sum w_\ell=1$ and
$\sum w_\ell\ell^{-j}=0$ for $1\le j\le m$.
Apply these identities to the derivative polynomial in
\eqref{eq:riesz-jet}. Its nonconstant terms cancel exactly.
Each remaining error is bounded by
$C_K(\ell T)^{D-m}\e^{-\sigma\ell T}$, and the finite weighted sum
has the bound in \eqref{eq:accelerated-bound}.
\end{proof}

The acceleration is algebraic and exact, not a fitted numerical
extrapolation. It uses $m+1$ cutoffs and no derivative values. Its
largest action is $(m+1)T$, so enumeration is linear in $T$ for fixed
$d,m$ before arithmetic precision is considered. An effective bit-cost
claim would additionally require certified representations of the
periods, reliable pole comparisons, and stable evaluation of the
retained finite sums. Those requirements are not supplied merely by
an analytic convergence theorem.

\section{Transport through a nonlinear inverse}
\label{sec:inverse}
The predecessor already provides an analytic inverse calculus after
block summation \cite{repo-block}. The result here is a quantitative
compatibility statement for the new raw-data regularization: one may
replace the block sum by an extrapolated finite Riesz sum and then
invert, without losing exponential accuracy.

\begin{theorem}[Certified local inverse stability]
\label{thm:inverse}
Fix a closed disk
$\overline{D}=\{z:|z-z_0|\le\rho\}$ and a larger concentric disk
$\overline{D^+}$ of radius $\rho+\eta$, where $\eta>0$ and
$\overline{D^+}\subset\Hh$. Put
\begin{equation}
 M=\sup_{\overline D}|\Ssum|,\qquad
 L=\sup_{\overline D}|\Ssum'|,
 \qquad |\varepsilon|M\le\rho/4,\quad
 |\varepsilon|L\le1/4.
 \label{eq:inverse-hypotheses}
\end{equation}
For each $|w-z_0|\le\rho/2$, the equation
\begin{equation}
 z+\varepsilon\Ssum(z)=w
 \label{eq:inverse-exact}
\end{equation}
has a unique solution in $\overline D$. For all sufficiently large $T$,
the equation
\begin{equation}
 z_T+\varepsilon\Accel_T^{(m)}(z_T)=w
 \label{eq:inverse-approximate}
\end{equation}
also has a unique solution there, where $m\ge d-1$. If
$\delta_T=\sup_{\overline{D^+}}|\Accel_T^{(m)}-\Ssum|$, then
\begin{equation}
 |z_T-z|\le
     \frac{|\varepsilon|}{1-|\varepsilon|L}\,\delta_T
   \le\frac{4|\varepsilon|}{3}\,\delta_T.
 \label{eq:inverse-error}
\end{equation}
In particular, the inverse error is
$O(|\varepsilon|T^{D-m}\e^{-\sigma T})$ uniformly in these targets,
where $\sigma=\min_{\overline{D^+}}\re z$.
\end{theorem}
\begin{proof}
The map $z\mapsto w-\varepsilon\Ssum(z)$ sends $\overline D$ into
its concentric disk of radius $3\rho/4$. On the convex disk its
Lipschitz constant is at most $|\varepsilon|L\le1/4$.
The contraction theorem gives the unique exact solution.

Corollary~\ref{cor:acceleration} makes $\delta_T\to0$. Cauchy's
estimate on circles of radius $\eta$ centered in $\overline D$ gives
\[
 \sup_{\overline D}|(\Accel_T^{(m)})'-\Ssum'|\le\delta_T/\eta.
\]
For large $T$, require
$|\varepsilon|\delta_T\le\rho/8$ and
$|\varepsilon|\delta_T/\eta\le1/4$. Then
$z\mapsto w-\varepsilon\Accel_T^{(m)}(z)$ maps $\overline D$ into
its disk of radius $7\rho/8$ and has Lipschitz constant at most $1/2$.
It too has a unique fixed point.

Subtract the two fixed-point equations and insert $\Ssum(z_T)$:
\[
 |z_T-z|\le|\varepsilon|L|z_T-z|
          +|\varepsilon|\delta_T.
\]
Solving this scalar inequality proves \eqref{eq:inverse-error}, and
\eqref{eq:accelerated-bound} gives the final rate. The argument also
shows, by the holomorphic implicit function theorem, that the local
inverse branches are holomorphic in the interior target disk.
\end{proof}

The factor $2\pi\ii$ converting $\Ssum$ into a Stokes jump can be
absorbed in $\varepsilon$. More general exactly invertible cores can
be treated in core coordinates, provided the image disks and derivative
bounds are supplied. That coordinate reduction and the abstract
contraction estimate are standard; the substantive new input is the
uniform exponentially accurate approximation in
Corollary~\ref{cor:acceleration}.

The theorem concerns a local analytic inverse. It does not assert global
univalence, a complete Borel singularity description for the inverse,
or convergence of its fully expanded raw action series. In particular,
when $\beta=\infty$, the unregularized linear residue contribution may
fail the term test, while the regularized finite inverse converges to
the correct local branch.

\section{Reproducible checks and numerical behavior}
\label{sec:verification}

\subsection{What was actually checked}
The supplied \path{verify.py} was executed with Python~3.13.5,
SymPy~1.14.0, and mpmath~1.3.0 at 110 decimal digits. It passed
\textbf{192 exact assertions} and \textbf{49 numerical assertions}.
The exact checks are distributed as follows:
\begin{center}
\begin{tabular}{lr}
\toprule
Finite identity family & Assertions\\
\midrule
Richardson moments, orders $1$ through $8$ &44\\
Continued-fraction determinant and complete-quotient identities &36\\
All-even parity pattern &31\\
Alternating parity pattern &30\\
Truncated-power collision scaling, $d=2,\ldots,7$ &18\\
Divided-difference polynomial reproduction &27\\
Weighted derivative identities &6\\
\midrule
Total &192\\
\bottomrule
\end{tabular}
\end{center}
The numerical checks comprise two reference-cutoff comparisons, ten
finite-bias remainders, ten extrapolated remainders, four cutoff
continuity diagnostics, fifteen actual sine-kernel collision scalings,
and eight finite Legendre classifications encountered up to $n=500$.
All tolerances and all measured values are recorded in
\path{results/verification.json}.

These are finite algebra checks and high-precision diagnostics, not
directed-rounding interval certificates. In particular, the program
does not approximate an entire rapidly growing continued-fraction
construction and infer its infinite exponent from a finite prefix.
The exponent and the endpoint types of those constructions are proved
in Section~\ref{sec:realization}. The program also does not turn the
analytic tail constants or the nonlinear inverse theorem into
machine-checked proofs.

\subsection{Derivative-bias removal}
Table~\ref{tab:riesz} uses $z=0.8$. For $d=2$, the periods are
$(1,\sqrt2)$, $D=2$, and $m=1$. For $d=3$, they are
$(1,\sqrt2,\sqrt3)$, $D=3$, and $m=2$. Define
\[
 Q_T(z)=\sum_{j=0}^m\binom mj\frac{\Ssum^{(j)}(z)}{T^j}.
\]
The reference sums and derivatives use poles below 260; comparisons
with a cutoff at 220 differed by less than $10^{-63}$ for the sum in
both models. This is a finite reference comparison, not a proof of a
$10^{-63}$ tail certificate. The displayed discrepancies are many
orders of magnitude larger than that comparison scale.

\begin{table}[htbp]
\centering\small
\begin{tabular}{rrrlll}
\toprule
$d$&$m$&$T$&$|\Riesz_T^{(m)}-\Ssum|$&$|\Riesz_T^{(m)}-Q_T|$&$|\Accel_T^{(m)}-\Ssum|$\\
\midrule
2 & 1 & 8 & $5.1888\times 10^{-4}$ & $3.9259\times 10^{-4}$ & $3.9543\times 10^{-4}$ \\
2 & 1 & 12 & $5.0987\times 10^{-5}$ & $3.3205\times 10^{-5}$ & $3.3214\times 10^{-5}$ \\
2 & 1 & 16 & $6.1721\times 10^{-5}$ & $1.4221\times 10^{-6}$ & $1.4221\times 10^{-6}$ \\
2 & 1 & 24 & $4.21\times 10^{-5}$ & $4.4517\times 10^{-9}$ & $4.4517\times 10^{-9}$ \\
2 & 1 & 32 & $3.1572\times 10^{-5}$ & $6.9309\times 10^{-12}$ & $6.9309\times 10^{-12}$ \\
3 & 2 & 8 & $4.538\times 10^{-6}$ & $2.0923\times 10^{-5}$ & $1.0766\times 10^{-5}$ \\
3 & 2 & 12 & $6.3028\times 10^{-6}$ & $4.3681\times 10^{-6}$ & $2.1823\times 10^{-6}$ \\
3 & 2 & 16 & $7.9841\times 10^{-6}$ & $7.5645\times 10^{-8}$ & $3.7821\times 10^{-8}$ \\
3 & 2 & 24 & $5.2087\times 10^{-6}$ & $4.3439\times 10^{-10}$ & $2.172\times 10^{-10}$ \\
3 & 2 & 32 & $3.8831\times 10^{-6}$ & $5.2882\times 10^{-13}$ & $2.6441\times 10^{-13}$ \\
\bottomrule
\end{tabular}

\caption{Finite Riesz bias and its extrapolated removal at $z=0.8$.
Ordinary high-precision arithmetic; no interval enclosure is implied.}
\label{tab:riesz}
\end{table}

The uncorrected error retains the algebraic derivative bias. Removing
that bias exposes the exponential scale predicted by
Theorem~\ref{thm:riesz}. There is no theorem that the error must decrease
monotonically at every small cutoff: a moving cutoff can cross or split
a block, and finite terms can cancel. The theorem is the uniform
upper estimate, not a pointwise monotonicity rule.

\subsection{A collision exactly at the cutoff}
Set $D=2$, $z=1.7$, $T=2$, and
\[
 \alpha_\pm(k)=\frac32\pm\sqrt2\,10^{-k}.
\]
The actions $2$ and $3/\alpha$ collide at the cutoff as $k\to\infty$.
Let $\Riesz_{2,\pm}^{(m)}$ denote the two finite sums. For $m=0$,
use the raw cutoff $P_2$. Table~\ref{tab:collision} illustrates all
three phenomena in Section~\ref{sec:sharpness}: the raw cutoff becomes
large; the first-order cutoff remains bounded but has different limits
from the two sides; the second-order difference tends to zero.

\begin{table}[htbp]
\centering\small
\begin{tabular}{r@{\qquad}llll}
\toprule
$k$&$P_{2,+}$&$\Riesz_{2,-}^{(1)}$&$\Riesz_{2,+}^{(1)}$&$|\Riesz_{2,+}^{(2)}-\Riesz_{2,-}^{(2)}|$\\
\midrule
4 & $47.805$ & $-9.352226\times 10^{-3}$ & $-4.842602\times 10^{-3}$ & $1.1465\times 10^{-6}$ \\
8 & $4.782\times 10^{5}$ & $-9.351778\times 10^{-3}$ & $-4.843219\times 10^{-3}$ & $1.1465\times 10^{-10}$ \\
16 & $4.782\times 10^{13}$ & $-9.351778\times 10^{-3}$ & $-4.843219\times 10^{-3}$ & $1.1465\times 10^{-18}$ \\
32 & $4.782\times 10^{29}$ & $-9.351778\times 10^{-3}$ & $-4.843219\times 10^{-3}$ & $1.1465\times 10^{-34}$ \\
\bottomrule
\end{tabular}

\caption{A cutoff collision distinguishes boundedness from continuity.}
\label{tab:collision}
\end{table}

At $\alpha=3/2$ itself, the sine kernel has a double pole at action two.
Using a simple-pole formula there would be invalid. The program avoids
that error by using irrational perturbations for the table. The
analytic discussion explains the correct multiple-pole residue and
why order one cannot have a side-independent cutoff limit.

\subsection{Rebuilding and the computational boundary}
The package includes the recorded table inputs, so the article can be
compiled without rerunning the computation. The commands
\begin{verbatim}
python verify.py --outdir build/results
pdflatex -interaction=nonstopmode -halt-on-error article.tex
pdflatex -interaction=nonstopmode -halt-on-error article.tex
pdflatex -interaction=nonstopmode -halt-on-error article.tex
\end{verbatim}
regenerate checks in a separate directory and typeset the recorded
article. The supplied build script directs TeX auxiliary output to
\path{build/}. No network calls or external data are used by the
verification program. Its simple-residue evaluator rejects a sine phase resolved as an
integer, including the tested exact coincidences; it is a diagnostic implementation,
not a general stable confluent-residue library.

A production numerical method should use the divided-difference or
contour representation for nearly colliding nodes, evaluate tail bounds
with outward rounding, and certify comparisons against the cutoff.
The present numerical tables should not be represented as having
performed those additional tasks.

\section{Further research questions and proposed projects}
\label{sec:questions}
The first question in the predecessor's list---ordinary convergence on
the critical line---has been answered above. The following projects
start at the remaining boundaries of the proved results rather than
restating that question unchanged.

\subsection*{1. The phase spectrum of the separate lattice sums}
For an explicitly prescribed amplitude sequence and parity pattern,
determine the set
\[
 E=\left\{y\in\R:
       \sum_k\eta_kA_k\e^{-\ii yq_k}\text{ converges}\right\}.
\]
Theorem~\ref{thm:classification} gives an exact reduction to this set,
not its measure, category, or Hausdorff dimension. A useful next result
would compare almost-everywhere and exceptional-point convergence for
superlacunary denominators while retaining the continued-fraction
realizability constraint.

\subsection*{2. Fixed-input minimal smoothing for three or more lattices}
Theorem~\ref{thm:sharp} proves the sharp order for estimates uniform in
the periods. For a fixed vector, smaller order may suffice. Seek a
necessary-and-sufficient criterion expressed through simultaneous
near-collisions and the effective multiplicities they create at large
action. This requires more than applying a two-lattice exponent
pairwise, because a cutoff can split a cluster of several poles.

\subsection*{3. Fractional cutoff orders and the optimal boundary profile}
The local scaling obstruction extends naturally to noninteger cutoff
orders, but the finite polynomial bias is special to integer $m$.
Develop a version for $(1-\lambda/T)_+^s$ with real $s$, giving its
sharp regularity threshold and a uniform bias expansion. Determine
which fractional combinations can improve error or reduce cancellation
without sacrificing confluence control.

\subsection*{4. Countably many lattices and unbounded cluster size}
For infinitely many period lattices the pigeonhole argument no longer
bounds cluster cardinality. Find a weighted criterion balancing the
number of nodes, regular-factor norms, and action growth. The aim is
a genuine countable-lattice theorem, not an invocation of
Theorem~\ref{thm:riesz} outside its fixed-$d$ hypothesis.

\subsection*{5. Algebraic and logarithmic Borel singularities}
Replace meromorphic poles by branch singularities and identify a cutoff
regularity rule that plays the role of multiplicity minus one. Such a
theorem must specify branches and continuation paths, replace residues
by Hankel contributions, and control merging singularities. It cannot
be obtained merely by substituting a noninteger pole order in the
present residue formulas.

\subsection*{6. Compatibility with transseries products and composition}
The local inverse theorem uses uniform approximation of one analytic
function. A broader algebraic objective is a summation map compatible
with products, differentiation, and composition on an explicit class
of raw action expansions. Identify the topology in which the smoothed
cutoffs define that map and determine when its actionwise coefficients
agree with formal operations. Local inverse stability alone does not
settle this algebraic compatibility.

\subsection*{7. Certified adaptive action budgets and bit complexity}
Make the constants in \eqref{eq:regular-bound} and
\eqref{eq:accelerated-bound} effective and computationally sharp for
specified period boxes. Combine interval representations of the periods,
certified cluster comparisons, confluent divided differences, and
outward-rounded tails. A useful output is an end-to-end algorithm whose
bit cost includes near-collision precision, not just the number of
retained actions.

\subsection*{8. The line of zero real part after regularization}
For the present $D\ge d$, ordinary raw convergence on $\re z=0$
already fails the term test. Investigate what boundary values, analytic
continuations, or distributional limits the Riesz sums possess there.
The exponential majorants used in this article disappear at that line;
a new oscillatory or distributional argument is required.

\subsection*{9. General multipliers and logarithmic coefficient growth}
A holomorphic multiplier with uniform tube bounds can replace $t^D$ in
parts of the block argument. Determine the endpoint extraction and
sharp Riesz tail for multipliers such as
$t^D(\log t)^r$ or regularly varying functions. Zeros of the multiplier
at convergent actions may alter the amplitude test, so a correct theorem
must include nonvanishing or explicit cancellation hypotheses.

\subsection*{10. A staged formalization in the repository}
The arithmetic and analytic halves can be developed independently.
First formalize the weighted denominator test and the primitive-pair
extraction. Separately formalize divided differences of truncated
powers with bounds that do not mention node separation. Then connect
the finite polynomial identity to block sums and package the inverse
stability estimate. The proof-dependency map below supplies concrete
milestones. No claim is made that these new theorems are already
available in Lean or that the required analytic interfaces have been
checked in the current Mathlib version.

\section{Proof audit, scope, and conclusions}
\label{sec:audit}

\subsection{Dependency map}
The main logical dependencies are short enough to record explicitly.
\begin{center}\small
\begin{tabular}{>{\raggedright\arraybackslash}p{.31\textwidth}>{\raggedright\arraybackslash}p{.61\textwidth}}
\toprule
Result & Essential inputs\\
\midrule
Critical absolute and term tests & Best approximation, exponential
weights, consecutive denominator intervals.\\[3pt]
Separate-series extraction & Legendre reduction, finite exponential
exponent, exact companion-pole identity, summable errors.\\[3pt]
Increasing-action classification & Two-node blocks, the individual term
test, the normally convergent block sum.\\[3pt]
Prescribed examples & Continued-fraction recurrence, bounded parity
rounding, a summable realization error.\\[3pt]
Riesz recovery & Bounded cluster size, pole-cancelled tube bounds,
Sobolev divided differences, a polynomial weighted tail.\\[3pt]
Sharpness and confluence & A first-pole collision family and the exact
regularity of the truncated power at its cutoff.\\[3pt]
Acceleration and inverse transport & Finite interpolation moments,
uniform approximation, Cauchy estimates, contraction.\\
\bottomrule
\end{tabular}
\end{center}

None of the main existence or convergence claims relies on a numerical
plot or on the observed behavior of a long but finite continued
fraction. Conversely, the exact finite identities checked by the
program do not verify all hypotheses of an infinite analytic theorem.
The final article is a conventional proof document, not a formal proof
assistant artifact.

\subsection{Claims not made}
The source comparison covers the named predecessor, its local
statements and research question, and the cited inventory. It is not
a claim that the entire repository was audited theorem by theorem.
The general theory of Riesz summability, the irrationality-base
formula, the divided-difference method, and the abstract local inverse
estimate are credited rather than relabeled as discoveries.

For fixed finite sine lattices, the new uniform Riesz theorem is
unconditional on Diophantine approximation. The endpoint classification,
however, concerns exactly two irrational lattices. No extension to
arbitrary countable singular supports, no full inverse-resurgence
theorem, and no global inverse univalence is asserted. The sharp
cutoff-order claim is explicitly about period-uniform gap-independent
control. Independent research priority and peer review remain separate
from the correctness of the written arguments.

\subsection{Conclusion}
The critical line is governed by a finer invariant than its exponential
abscissa. Continued fractions supply that invariant in the sequence
$q_k^Dq_{k+1}\e^{-bq_k}$. Vanishing of these amplitudes controls unfinished
pairs; summability controls absolute convergence; their parity and
phases control the two separate lattice sums. This separates three
questions that the open-half-plane root test cannot distinguish.

The cutoff theorem provides a complementary answer that does not ask
for small-divisor arithmetic at all. A truncated power with enough
regularity controls the divided differences of a block even when the
cutoff splits it. Its finite derivative bias can be cancelled exactly,
yielding an exponentially accurate finite-action representation and a
stable nonlinear inverse approximation. In this setting, both the
arithmetic boundary and the analytic regularization can therefore be
made precise without identifying a formal transseries with an
unconditionally convergent raw numerical sum.

\appendix
\section{Repository and primary-literature provenance}
\label{app:sources}
The pinned repository commit is
\begin{center}
\texttt{3d5973524506411392a911470b5ddc35521568ea}.
\end{center}
The directly inspected predecessor source is
\begin{quote}\small
\path{Analysis/Transseries/docs/series-and-transseries/Resonance_Block_Summation_Transseries/article.tex}.
\end{quote}
Its Git blob identifier is
\texttt{934422729a618bf2bb6792f4341ec73b9f5cf94e}.
The inspected ranges include lines 1--240 (scope and conventions),
440--710 (divided differences, block summation, residues, abscissa),
1110--1370 (inverse transport and the general bounded-cluster criterion),
and 1371--1580 (stable evaluation and the explicit boundary research
question). Its package README and the group inventory were also read
through the connected GitHub tools. The repository was not modified.

For the external primary sources, Sondow's continued-fraction formula
and Legendre lemma were checked in the PDF, including the displayed
formula for the logarithm of the irrationality base. De Boor's
Genocchi--Hermite formula and the contour/divided-difference discussion
were inspected. The Defant--Schoolmann paper was used for the established
Riesz-summability context, not as a black-box proof of
Theorem~\ref{thm:riesz}. These references establish antecedents and
notation; a comprehensive priority review of all specialized
small-divisor and summability literature has not been performed.

\begin{thebibliography}{9}

\bibitem{repo-block}
\repo\ repository,
\emph{Resonance-Block Summation of Transseries: Small divisors,
coalescing Stokes actions, and nonlinear inversion},
research report filed 29 September 2026, with editorial amendments.
Source pinned at commit \texttt{3d5973524506411392a911470b5ddc35521568ea}.
Especially \texttt{thm:block}, \texttt{thm:arithmetic},
\texttt{thm:inverse}, and further-research Question~1.
\url{https://github.com/VladimirReshetnikov/ProveIt/blob/3d5973524506411392a911470b5ddc35521568ea/Analysis/Transseries/docs/series-and-transseries/Resonance_Block_Summation_Transseries/article.tex}.

\bibitem{repo-index}
\repo\ repository,
\emph{Series and transseries}, group inventory,
\path{Analysis/Transseries/docs/series-and-transseries/README.md},
same pinned commit, inspected 29 September 2026.
\url{https://github.com/VladimirReshetnikov/ProveIt/blob/3d5973524506411392a911470b5ddc35521568ea/Analysis/Transseries/docs/series-and-transseries/README.md}.

\bibitem{sondow}
J.~Sondow,
\emph{Irrationality Measures, Irrationality Bases, and a Theorem of
Jarn\'ik}, arXiv:math/0406300 (2004).
See Lemma~1 and Theorem~1, especially equation~(3).
\url{https://arxiv.org/abs/math/0406300}.

\bibitem{deboor}
C.~de~Boor,
\emph{Divided Differences},
Surveys in Approximation Theory \textbf{1} (2005), 46--69.
See the contour representation and the Genocchi--Hermite formula~(52).
\url{https://arxiv.org/abs/math/0502036}.

\bibitem{defant-schoolmann}
A.~Defant and I.~Schoolmann,
\emph{Riesz Summability on Boundary Lines of Holomorphic Functions of
Finite Order Generated by Dirichlet Series},
arXiv:2107.10145, version~3 (2022).
The quoted title is the title printed in the inspected PDF.
\url{https://arxiv.org/abs/2107.10145}.

\end{thebibliography}
\end{document}
